\chapter{El sonido: velocidad, tono e intensidad sonora}\label{ch:g8:sound-pitch-loudness}

Las cuerdas cortas cantan agudo y los pellizcos fuertes cantan alto ---
tu guitarra de gomas elásticas encontró las reglas hace años. Este año
las reglas reciben números: una cuenta por segundo para el tono, un
tamaño de vaivén para el volumen y una ventana --- con bordes --- sobre
lo que los oídos humanos pueden captar. Más allá de uno de los bordes,
los murciélagos se ríen de nosotros.

\section{La frecuencia: el tono, contado}

\begin{definition}[Frecuencia]\label{def:g8:sound-pitch-loudness:frequency}
La \emph{frecuencia}\index{frecuencia} $f$ de una \omterm{def:g3:sound-around-us:vibration}{vibración} es
el número de vaivenes completos que hace cada segundo. Su
\omterm{def:g6:measurement-in-science:measuring}{unidad} es el \emph{hercio}\index{hercio} (\unit{Hz}): un vaivén
por segundo --- en honor del físico que fabricó y capturó las primeras
ondas invisibles. Una cuerda de guitarra que tiembla $110$ veces por
segundo vibra a \qty{110}{Hz}; las alas de un mosquito, cerca de
\qty{600}{Hz}; el aleteo lento de una bandera grande, unos pocos
hercios.
\end{definition}

\begin{proposition}[El tono es la frecuencia]\label{prop:g8:sound-pitch-loudness:pitch}
Lo que el oído llama \emph{tono} es la \omterm{def:g8:sound-pitch-loudness:frequency}{frecuencia} y nada más:
cuanto más rápido tiembla la fuente, más aguda es la nota. Las reglas
de las gomas elásticas se traducen al instante --- lo corto, lo tenso y
lo delgado tiemblan con $f$ más alta. El ancla de la música: los
instrumentos de concierto del mundo entero se afinan con el la de
\qty{440}{Hz}. Duplicar la \omterm{def:g8:sound-pitch-loudness:frequency}{frecuencia} de cualquier nota la sube el
más consonante de todos los intervalos --- la octava: otra vez un la,
una octava más arriba, a \qty{880}{Hz}.
\end{proposition}

\begin{omfigure}
\begin{tikzpicture}
\begin{axis}[omaxis, width=10.8cm, height=4.4cm,
    xlabel={tiempo (\unit{ms})}, ylabel={vaivén},
    xtick={0,5,10,15,20}, ytick=\empty,
    ymin=-1.5, ymax=1.5, xmin=0, xmax=21, samples=200]
  \addplot[very thick, omDef, domain=0:20]
    {sin(deg(2*3.14159*0.1*x))};
  \addplot[very thick, omThm, domain=0:20]
    {0.9*sin(deg(2*3.14159*0.4*x))};
\end{axis}
\end{tikzpicture}

{\small Dos \omterm{def:g3:sound-around-us:sound}{sonidos} dibujados por un micrófono: en los mismos $20$
milisegundos, el trazo azul oscila dos veces (\qty{100}{Hz}, un
zumbido grave) mientras el rojo oscila ocho (\qty{400}{Hz}, una nota
cantada). El tono es la cuenta.}
\end{omfigure}

\begin{method}[Ver el sonido]\label{met:g8:sound-pitch-loudness:seeing}
El tambor con granos de arroz enseñaba el temblor del \omterm{def:g3:sound-around-us:sound}{sonido}; un
micrófono lo dibuja.
\begin{enumerate}
  \item pon en marcha una aplicación de análisis de \omterm{def:g3:sound-around-us:sound}{sonido} en un
        móvil o en un ordenador --- su pantalla representa el vaivén
        del micrófono frente al tiempo, un osciloscopio moderno;
  \item tararea grave y sostenido: un trazo perezoso y muy espaciado;
        canta agudo: los picos se apretujan;
  \item lee el indicador de \omterm{def:g8:sound-pitch-loudness:frequency}{frecuencia} mientras silbas: casi
        todos los silbidos aterrizan entre $500$ y \qty{2000}{Hz};
  \item golpea el diapasón marcado ``la 440'' y comprueba que el
        indicador honra su grabado.
\end{enumerate}
Una sola pantalla sustituye a un capítulo de conjeturas: el tono
contado ante tus ojos.
\end{method}

\section{La ventana del oído}

\begin{proposition}[El intervalo de la audición humana]\label{prop:g8:sound-pitch-loudness:range}
Los oídos jóvenes y sanos captan vibraciones de entre unos
\qty{20}{Hz} y \qty{20000}{Hz}. Por debajo de la ventana están los
\emph{infrasonidos}\index{infrasonido} --- notados a veces como un
retumbar en el pecho, oídos nunca; por encima, los
\emph{ultrasonidos}\index{ultrasonido} --- silencio para nosotros, por
violento que sea el temblor. La ventana se estrecha con la edad y con
el maltrato: la octava más aguda suele ser la primera víctima de las
costumbres ruidosas.
\end{proposition}

\begin{example}[Más allá de los bordes]\label{ex:g8:sound-pitch-loudness:beyond}
La naturaleza y la industria trabajan las dos más allá de nuestra
ventana. Los murciélagos cazan con chasquidos ultrasónicos cercanos a
los \qty{50000}{Hz} y leen \omterm{ex:g4:how-sound-travels:echo}{ecos} más finos de lo que podría pintar
ningún \omterm{def:g3:sound-around-us:sound}{sonido} audible --- y los delfines igual, con el espíritu
del capítulo del sonar. Los perros responden a los silbatos
``silenciosos'' afinados justo por encima de nuestro borde. Los
escáneres médicos hacen imágenes de los niños por nacer con
\omterm{ex:g4:how-sound-travels:echo}{ecos} inofensivos de \omterm{prop:g8:sound-pitch-loudness:range}{ultrasonidos}. En el otro borde, los
elefantes conversan con retumbos infrasónicos a lo largo de kilómetros,
y las tormentas y los terremotos se anuncian en tonos a los que no
atiende ningún oído.
\end{example}

\section{El volumen: el tamaño del vaivén}

\begin{proposition}[El volumen es la amplitud]\label{prop:g8:sound-pitch-loudness:loudness}
Lo que el oído llama \emph{\omterm{def:g6:mass-and-volume:volume}{volumen}} corresponde a la
\emph{amplitud}\index{amplitud} de la \omterm{def:g3:sound-around-us:vibration}{vibración} --- el tamaño del
vaivén, la \omterm{def:g2:pushes-and-pulls:force}{fuerza} del apretujamiento viajero. Pellizca con suavidad o
con saña: los picos del trazo quedan bajos o altos mientras su
\emph{espaciado} --- el tono --- se queda intacto. La
\omterm{prop:g8:sound-pitch-loudness:loudness}{amplitud} y la \omterm{def:g8:sound-pitch-loudness:frequency}{frecuencia} son mandos independientes:
cualquier tono se puede susurrar o berrear.
\end{proposition}

\begin{omfigure}
\begin{tikzpicture}
\begin{axis}[omaxis, width=10.8cm, height=4.4cm,
    xlabel={tiempo (\unit{ms})}, ylabel={vaivén},
    xtick={0,5,10,15,20}, ytick=\empty,
    ymin=-1.5, ymax=1.5, xmin=0, xmax=21, samples=200]
  \addplot[very thick, omDef, domain=0:20]
    {0.35*sin(deg(2*3.14159*0.25*x))};
  \addplot[very thick, omThm, domain=0:20]
    {1.2*sin(deg(2*3.14159*0.25*x))};
\end{axis}
\end{tikzpicture}

{\small La misma nota, susurrada y berreada: picos igual de espaciados
(un solo tono, \qty{250}{Hz}) y de alturas desiguales --- el volumen es
el tamaño del vaivén.}
\end{omfigure}

\begin{example}[La escalera de los decibelios]\label{ex:g8:sound-pitch-loudness:decibels}
El \omterm{def:g6:mass-and-volume:volume}{volumen} se gradúa en \emph{decibelios}\index{decibelio} (\unit{dB})
--- una escalera construida para el asombroso intervalo del oído, en la
que cada refuerzo del \omterm{def:g3:sound-around-us:sound}{sonido} por un factor de diez añade el mismo
número de peldaños. Referencias: hojas que susurran, \qty{10}{dB}; una
conversación tranquila, \qty{50}{dB}; tráfico intenso, \qty{80}{dB};
las primeras filas de un concierto de rock, \qty{110}{dB}; el dolor y
el daño instantáneo, cerca de \qty{120}{dB}. La extraña aritmética de
la escalera --- pasos iguales para \omterm{def:g2:pushes-and-pulls:force}{fuerzas} diez veces mayores --- es un
logaritmo, uno de los grandes inventos que te esperan en las
matemáticas del bachillerato.
\end{example}

\begin{remark}[El presupuesto del oído]\label{rem:g8:sound-pitch-loudness:earsafety}
El tamborcito de tu oído, que conociste en el capítulo de tercero,
lleva las cuentas de toda una vida. Por encima de unos \qty{85}{dB}, el
daño se acumula con el tiempo de exposición --- el límite de una jornada
en una fábrica ruidosa, solo minutos a los niveles de la primera fila
de un concierto; los auriculares a todo volumen se sitúan de lleno en
la banda de peligro. El oído perdido por los vaivenes altos no
vuelve nunca: las notas más agudas se van primero, en silencio y para
siempre. Los oídos no tienen párpados --- ni piezas de repuesto.
\end{remark}

\begin{example}[La velocidad, a la que no toca ningún mando]\label{ex:g8:sound-pitch-loudness:speed}
Ninguno de los dos mandos toca la entrega. Agudos y graves, fuertes y
suaves, todos los \omterm{def:g3:sound-around-us:sound}{sonidos} viajan en el único relevo del
medio al único paso del medio --- \qty{340}{m/s} en el
\omterm{def:g2:air-around-us:air}{aire}. La demostración se representa cada noche: una orquesta oída
al otro lado de un parque llega \emph{afinada y a tiempo} --- si las
notas agudas fueran más veloces que las graves, todos los acordes
lejanos se emborronarían en arpegios. (El \omterm{def:g6:mass-and-volume:volume}{volumen} se apaga con la
distancia --- las cáscaras que se extienden del capítulo de séptimo ---,
pero apagarse es debilitarse, nunca frenar.)
\end{example}

\section{Ejercicios}

\begin{exercise}[$\star$]\label{exo:g8:sound-pitch-loudness:1}
Define la \omterm{def:g8:sound-pitch-loudness:frequency}{frecuencia} y su \omterm{def:g6:measurement-in-science:measuring}{unidad}. ¿Cuál es el tono del la
de concierto y el del la de una octava por encima?
\end{exercise}

\begin{exercise}[$\star$]\label{exo:g8:sound-pitch-loudness:2}
Dos trazos abarcan los mismos $20$ milisegundos: uno enseña $4$
vaivenes completos y el otro $16$. Calcula las dos \omterm{def:g8:sound-pitch-loudness:frequency}{frecuencias}; ¿cuál
suena más agudo?
\end{exercise}

\begin{exercise}[$\star$]\label{exo:g8:sound-pitch-loudness:3}
Da la ventana de la audición humana y los nombres de los dos países
que hay más allá de sus fronteras --- con un habitante de cada uno.
\end{exercise}

\begin{exercise}[$\star$]\label{exo:g8:sound-pitch-loudness:4}
En la pantalla del analizador, ¿qué cambia cuando cantas la misma nota
más fuerte? ¿Y cuando cantas una nota más aguda con el mismo \omterm{def:g6:mass-and-volume:volume}{volumen}?
\end{exercise}

\begin{exercise}[$\star$]\label{exo:g8:sound-pitch-loudness:5}
Traduce al lenguaje de la \omterm{def:g8:sound-pitch-loudness:frequency}{frecuencia} las viejas reglas de las gomas
elásticas: ¿qué le hacen a $f$ lo corto, lo tenso y lo delgado?
\end{exercise}

\begin{exercise}[$\star$]\label{exo:g8:sound-pitch-loudness:6}
Sitúa en la escalera de los \omterm{ex:g8:sound-pitch-loudness:decibels}{decibelios}: las hojas que susurran; el
tráfico; la primera fila de un concierto; el dolor. ¿Dónde empieza la
banda de peligro?
\end{exercise}

\begin{exercise}[$\star$]\label{exo:g8:sound-pitch-loudness:7}
¿Por qué a sus dueños les parecen estropeados los silbatos para perros?
¿Y por qué el perro no está de acuerdo?
\end{exercise}

\begin{exercise}[$\star$]\label{exo:g8:sound-pitch-loudness:8}
Una orquesta lejana llega afinada y a tiempo. ¿Qué demuestra esto sobre
la \omterm{def:g7:motion-average-speed:formula}{velocidad} del \omterm{def:g3:sound-around-us:sound}{sonido} para los distintos tonos y
volúmenes?
\end{exercise}

\begin{exercise}[$\star\star$]\label{exo:g8:sound-pitch-loudness:9}
Un murciélago chasquea a \qty{50000}{Hz}. ¿Cuántos vaivenes completos
contiene un chasquido de \qty{0.01}{s}? ¿Por qué tiene que usar un
temblor tan rápido la formación de imágenes finas por \omterm{ex:g4:how-sound-travels:echo}{eco}?
(Acuérdate de qué tienen en común los pinceles más finos.)
\end{exercise}

\begin{exercise}[$\star\star$]\label{exo:g8:sound-pitch-loudness:10}
La tecla más grave de un piano de cola vibra cerca de \qty{27}{Hz}; la
más aguda, cerca de \qty{4200}{Hz}. Sitúa las dos dentro (o frente a)
la ventana de la audición --- y explica por qué un afinador de
pianos que envejece puede necesitar antes la ayuda del analizador en un
extremo del teclado.
\end{exercise}

\begin{exercise}[$\star\star$]\label{exo:g8:sound-pitch-loudness:11}
Las normas de los conciertos permiten $8$ \omterm{ex:g2:measuring-time:clock}{horas} a \qty{85}{dB}
pero solo minutos cerca de \qty{110}{dB}, y el presupuesto se divide
por dos cada pocos peldaños. ¿Qué clase de crecimiento se esconde en
los pasos iguales de la escalera --- y qué enseña el derrumbe del
presupuesto sobre eso de ``solo un pelín más alto''?
\end{exercise}

\begin{exercise}[$\star\star\star$]\label{exo:g8:sound-pitch-loudness:12}
Diseña la prueba de \omterm{def:g3:sound-around-us:sound}{sonido} completa del salón de actos de un colegio
con un único móvil con analizador: cómo verificar el la del piano
frente a los $440$; cómo hallar el volumen del salón en la última fila
durante una prueba de batería; cómo pillar si la ventilación emite un
retumbo infrasónico que el oído no capta; y cuál de los tres hallazgos
está \emph{menos} equipado para certificar el micrófono del móvil
(consulta la propia ventana de audición del móvil --- construida
para voces).
\end{exercise}

\section{Problema: la prueba de sonido}

\begin{problem}[{Problema de fin de semana --- prueba de sonido en el
festival del colegio; un analizador, tres grupos y una enfermera
nerviosa; la velada de la mesa de mezclas}]\label{pb:g8:sound-pitch-loudness:1}
Víspera del festival: tú manejas el analizador en la mesa de mezclas,
con la enfermera del colegio vigilando el medidor de \omterm{ex:g8:sound-pitch-loudness:decibels}{decibelios} por
encima de tu hombro.

\medskip\noindent\textbf{Parte I --- Hora de afinar.}
\begin{enumerate}
  \item El diapasón del coro canta a \qty{440}{Hz}; el analizador
        enseña el la del piano a \qty{432}{Hz}. Diagnostica el piano
        y nombra el mando --- tono o \omterm{def:g6:mass-and-volume:volume}{volumen} --- que necesita al
        afinador.
  \item La cuerda más grave del bajo marca \qty{41}{Hz}. Sitúala en
        la ventana de la audición --- y explica por qué el
        público la \emph{notará} en el pecho tanto como la oirá.
  \item La nota sostenida del cantante enseña picos espaciados
        exactamente como los del diapasón, pero tres veces más
        altos. Compara los dos \omterm{def:g3:sound-around-us:sound}{sonidos} en los dos mandos.
  \item Un técnico joven afirma que los altavoces nuevos del salón
        ``reproducen hasta \qty{22000}{Hz}''. ¿Quién del público
        podrá certificar la parte alta de ese intervalo --- y quién
        no?
\end{enumerate}

\medskip\noindent\textbf{Parte II --- El libro de cuentas de la
enfermera.}
\begin{enumerate}[resume]
  \item Niveles del ensayo: sesión acústica \qty{75}{dB}; sesión de
        rock \qty{100}{dB}; final previsto ``un pelín por encima'',
        a \qty{110}{dB}. ¿Qué sesiones marca la enfermera, y contra
        qué línea del presupuesto?
  \item La enfermera reparte tapones de espuma que atenúan los
        \omterm{def:g3:sound-around-us:sound}{sonidos} \qty{20}{dB}. ¿En qué se convierte la sesión de
        rock detrás de ellos --- y por qué los técnicos sensatos
        llevan tapones en todas las actuaciones?
  \item Al ponerse al doble de distancia del escenario, el medidor
        baja de forma perceptible. ¿Qué vieja ley de las cáscaras
        que se extienden está actuando --- y qué mando (tono o
        \omterm{def:g6:mass-and-volume:volume}{volumen}) gira la distancia?
  \item El batería protesta: ``¡Más flojo significa más apagado ---
        los agudos mueren primero cuando tocamos suave!''. Desenrédale
        los dos mandos: ¿qué cambia en realidad tocar suave?
\end{enumerate}

\medskip\noindent\textbf{Parte III --- Física de festival.}
\begin{enumerate}[resume]
  \item Los fuegos artificiales cierran el espectáculo a
        \qty{680}{m}. Calcula el retraso del destello al estallido
        que notará la multitud --- y la \omterm{def:g8:sound-pitch-loudness:frequency}{frecuencia} del retumbo
        si los cohetes grandes golpean dos veces por segundo
        (comprobación de ventana: ¿se oye o se nota?).
  \item El último misterio de la noche: un vídeo del concierto
        grabado con el móvil suena flaco al reproducirlo --- ha
        desaparecido el retumbo del pecho. Consulta la advertencia
        del ejercicio 12: ¿qué es lo que el micrófono no se
        construyó nunca para captar?
  \item El informe de la enfermera pide ``una frase para cada
        cosa'' sobre los dos mandos de la velada y sobre la tercera
        cosa que ninguno de los dos mandos puede tocar. Redáctalas.
  \item Recogida: el ahuyentador ultrasónico de plagas del
        conserje, en marcha toda la tarde junto a la puerta del
        escenario, no molestó a ningún oyente y, según dicen, a
        muchos ratones. Explica las dos audiencias con una sola
        proposición sobre la ventana.
\end{enumerate}
\end{problem}
